\documentclass[12pt, a4paper, titlepage, abstract=on]{scrreprt}

\usepackage[utf8]{inputenc}
\usepackage[czech]{babel}
\usepackage[T1]{fontenc}
\usepackage{microtype}
\usepackage{cmbright}
\usepackage{csquotes}
\usepackage{geometry}
\usepackage{graphicx}
\usepackage{hyperref}
\usepackage{tabularx}
\usepackage{minted}
\usepackage{tikz}
\geometry{margin=5em}
\usetikzlibrary{positioning}
\usetikzlibrary{matrix}
\setminted{
    fontsize=\small
}
\title{Compiler vlastního programovacího jazyka}
\author{Ondřej Suk}
\date{2026}

\begin{document}

\maketitle

\begin{abstract}
Práce se zabývá vytvářením jednoduchého compileru programovacího jazyka vytvořeného specificky pro tento projekt za pomocí knihovny LLVM.
Jazyk samotný je velmi jednoduchý a je velmi obtížné v něm programovat složitější programy. Důvodem je, že i v aktuálním stavu je projekt poměrně
rozsáhlý a tudíž obsahuje jenom minimum, aby v něm bylo možné programovat. Některé části jazyka jsou navrženy netradičním způsobem, což je z důvodu, že na projektu plánuji pracovat i po maturitní zkoušce, tudíž jsem kód přizpůsobil tomu, jak plánuji, aby fungoval finální jazyk.

Compiler jsem implementoval v jazyce C++ 17. Vstupem compileru je soubor s kódem. Ten je potom přeložen do LLVM IR, jazyka poskytovaného knihovnou
LLVM, ze kterého je pak zkompilován na objektový kód. Z objektového kódu nakonec vznikne spustitelný soubor. Kompilace z LLVM IR na binární kód je
již součástí knihovny LLVM, tudíž se hlavní část projektu zabývá překladem souboru s programem na LLVM IR.
\end{abstract}

\tableofcontents
\clearpage

\chapter{Úvod}
Cílem projektu bylo vytvořit compiler jednoduchého jazyka, který bude sloužit jako základ mnohem pokročilejšího jazyka, na kterém plánuji
pracovat po maturitě. Z toho důvodu je celý projekt přizpůsobený mé trochu netradiční vizi pro syntax finálního jazyka. Při tvorbě projektu
jsem chtěl objevovat nové koncepty sám, takže se výsledek může lišit od běžných compilerů. V jazyce tedy není v aktuálním stavu snadné programovat
a mnoho věcí, které by programátor očekával, chybí. To je v souladu s cílem projektu, protože plně rozsáhlý programovací jazyk by potřeboval
množství času mnohem vyšší, než je v rámci maturitního projektu možné.

Již zpočátku bylo jasné, že použiji knihovnu LLVM. LLVM slouží jako základ mnoha moderních compilerů. Poskytuje jazyk zvaný LLVM IR, do kterého
compiler přeloží vstupní kód. LLVM pak může kód zkompilovat pro kteroukoli cílovou platformu, takže je práce pro podporu mnoha platforem snazší.
Navíc dokáže LLVM vzniklý kód optimalizovat. Jako programovací jazyk jsem si vybral C++, protože LLVM je primárně knihovnou pro jazyk C++.
Specificky je LLVM napsáno pro C++ 17, a tak jsem i já použil C++ 17. Běh programu jsem strukturoval na několik fází a kostru, která ty fáze spojuje.

Hotový projekt tedy napřed přečte soubor s kódem a postupně během několika fází z něj extrahuje všechny důležité informace. Na základě těch informací
pak napíše program přeložený do jazyka LLVM IR. Poslední fáze potom zavolá knihovnu LLVM a ta z LLVM IR vytvoří binární kód specifický pro cílovou
platformu. Následně pak z binárního kódu vytvoří spustitelný soubor. (Na Windows soubor s příponou .exe)

Hlavní část dokumentace je rozdělena na dvě části. Programátorská dokumentace se zabývá samotným compilerem. Popisuje detailně jednotlivé fáze
kompilace, a proč jsem některé věci navrhl určitým způsobem. Měla by poskytovat vše důležité pro další práci na projektu. Uživatelská dokumentace
se zaměřuje na samotný jazyk a vše, co obsahuje.
\chapter{Programátorská dokumentace}
\section{Použité technologie nástroje}
Níže je seznam použitých nástrojů. Je členěn na tři části. První oddíl obsahuje nástroje, které mají v projektu viditelnou stopu. Druhý je věnován nástrojům,
kterých jsem využil při programování na Linuxu, ale místo nich je k editaci projektu snadné dosadit nástroje jiné. Poslední oddíl se věnuje mému přístupu k
umělé inteligenci při práci na projektu.
\subsection{Uvnitř projektu}
\paragraph{\href{https://llvm.org/}{LLVM}}
LLVM je knihovna jazyka C++, která slouží jako základ compileru. Definuje jazyk LLVM IR, který je nízkoúrovňovou, ale zároveň platformově nezávislou
reprezentaci kompilovaného kódu.
\paragraph{\href{https://lld.llvm.org/}{LLD}}
LLD je linker, který je součástí projektu LLVM. Linker zpracovává soubory s binárními instrukcemi a vytváří z nich spustitelné soubory. (.exe soubory)
\paragraph{\href{https://github.com/gouwsxander/easy-args}{Easyargs}}
Knihovnu Easyargs jsem použil hlavně, protože poskytuje nenáročné rozhraní pro definování argumentů příkazového řádku. Uživatelské rozhraní nepovažuji
za důležitou část projektu, tudíž jsem vybral jednoduchou knihovnu.
\paragraph{\href{https://www.gnu.org/software/make/manual/make.html}{Makefile}}
Makefile je systém, pomocí kterého jsem vytvořil kompilační skript pro tento projekt. Vybral jsem si ho, protože na rozdíl od jiných C++ build systémů
umožňuje Makefile specifikovat přesně ty příkazy, které budou při kompilaci projektu použity.
\paragraph{\href{https://www.mingw-w64.org/}{MinGW}}
MinGW je port nástrojů příkazového řádku Linuxu pro Windows. Podrobnosti, jak jsem MinGW použil jsou v pozdější části dokumentace.
\paragraph{\href{https://clang.llvm.org/}{Clang}}
Jako C++ compiler jsem si vybral Clang, protože je součástí LLVM, a protože poskytuje nástroje Clang Format a Clang Tidy, které jsem také použil.
\paragraph{\href{https://clang.llvm.org/docs/ClangFormat.html}{Clang Format}}
Clang Format je nástroj, díky kterému zajišťuje, že všechen kód v projektu je správně zformátovaný.
\paragraph{\href{https://clang.llvm.org/extra/clang-tidy/}{Clang Tidy}}
Clang Tidy je nástroj, který pomáhá s hledáním chyb v kódu ještě před tím, než je vůbec zkompilován.
\paragraph{\href{https://git-scm.com/}{Git}}
Pro zálohování změn v projektu jsem využil Git a projekt je veřejně viditelný přes repozitář na GitHubu: \href{https://github.com/F100cTomas/Lang}{https://github.com/F100cTomas/Lang}.
\paragraph{\href{https://www.latex-project.org/}{\LaTeX}}
\LaTeX je můj primární nástroj pro psaní dokumentace.
\subsection{Pro vývoj na Linuxu}
\paragraph{\href{https://neovim.io/}{Neovim}}
Neovim sloužil jako můj primární editor pro psaní kódu, hlavně tedy C++.
\paragraph{\href{https://apps.kde.org/kate/}{Kate}}
Kate sloužil jako můj sekundární editor pro psaní kódu. Psal jsem v něm hlavně kompilační skript a testovací programy. Vybral jsem si ho, protože
je výchozím editorem pro KDE desktop.
\paragraph{\href{https://apps.kde.org/kile/}{Kile}}
Kile sloužil jako můj editor pro psaní dokumentace. Vybral jsem si ho, protože je výchozím \LaTeX\ editorem pro KDE desktop.
\paragraph{\href{https://lldb.llvm.org/}{LLDB}}
LLDB je debugger, který jsem použil jako poslední možnost při diagnóze chyb v kódu, převážně segmentačních chyb. Vybral jsem si ho, protože je součástí LLVM.
\paragraph{\href{https://archlinux.org/}{Arch Linux}}
Arch Linux je Linux distribuce, na které probíhala většina vývoje. Použil jsem ho, protože je již mým primárním operačním systémem.
\subsection{Pro vývoj na Windows}
\paragraph{\href{https://www.qemu.org/}{QEMU}}
Protože nemám doma dostupný počítač s instalací Windows, probíhal téměř veškerý vývoj pro Windows uvnitř virtuálního prostředí Windows 11 v QEMU. Je nutné
podotknout, že to nebyla ideální volba, která vedla k častému zamrzání obrazovky a dalším nepříjemnostem.
\paragraph{\href{https://notepad-plus-plus.org/}{Notepad++}}
Notepad++ jsem si vybral, protože se jedná o dobře optimalizovaný editor. Nechtěl jsem totiž své již tak omezené Windows prostředí zbytečně zatěžovat.
\subsection{Ohledně AI}
Přestože projekt neobsahuje jediný řádek C++ napsaný AI, měla umělá inteligence v mém projektu zásadní roli. Jako poskytovatele AI jsem si vybral ChatGPT,
protože jsem ChatGPT používal už před začátkem projektu. Hlavní roli měla AI v debugování kódu. Umělá inteligence má totiž nadlidskou schopnost hledat hrubky.
Dále mi umělá inteligence pomohla pochopit věci, které jsem se v rámci projektu učil nově. AI se hodně podílela na sepsání Makefile.
\section{Základní struktura}
Projekt je rozdělen na dvě hlavní části. První je kostra, která definuje základní běh programu, a druhá jsou jednotlivé fáze, které implementují
jednoduché rozhraní pro konverzi kódu z jedné reprezentace na druhou. Tyto reprezentace jsou si docela blízké, ale jejich vrstvením dokážou
zdrojový kód přetvořit na velice rozdílnou formu, spustitelný soubor.

Ve složce s compilerem se nachází několik důležitých souborů. Prvním je \texttt{Makefile}, který obsahuje skript, pomocí kterého je možné
celý projekt zkompilovat. Dále se zde nachází \texttt{.clang-format} a \texttt{.clang-tidy}, které obsahují konfigurace pro nástroje
používané při vývoji. Méně důležitým souborem je \texttt{kernel32.def}, který obsahuje seznam všech Windows API funkcí, které může zkompilovaný
projekt zavolat. \texttt{kernel32.def} existuje pouze kvůli kompatibilitě s Windows. V poslední řadě tu je soubor \texttt{code}, který by
vždy měl obsahovat příkladový program k testování projektu.

Všechen zdrojový kód compileru se nachází ve složce \texttt{src/}. Soubory, které jsou součástí kostry, a tudíž nepatří žádné kompilační
fázi jsou přímo ve složce \texttt{src/} a každá kompilační fáze má vlastní podsložku uvnitř \texttt{src/}. Na rozdíl od některých C++
projektů se hlavičkové soubory nachází ve stejné složce, jako ostatní soubory s kódem. Soubory s C++ kódem mají přípony \texttt{.cpp} a
\texttt{.hpp}.

\vspace{1em}
\noindent Následuje souhrn kompilačních fází se jmény, která mají v kódu, a jejich krátké vysvětlení:
\begin{itemize}
 \item Lexer -- Rozsekává vstupní text na slova. (tokeny)
 \item Preparser -- Parsuje syntax klíčových slov a sestavuje tabulku symbolů.
 \item Parser -- Řeší pořadí operátorů a vytváří abstraktní syntaktický strom. (AST)
 \item Code Generator -- Vytváří instrukce LLVM IR na základě AST.
 \item LLVM Interface -- Zkompiluje vzniklé LLVM IR do spustitelného souboru.
\end{itemize}

\begin{figure}[H]
\centering
\begin{tikzpicture}[
    node distance=2mm,
    every node/.style={draw, rectangle, rounded corners, minimum width=3mm, minimum height=12mm, align=center},
    end/.style={draw, circle, minimum width=16mm, minimum height=1mm, align=center},
    arrow/.style={->, thick}
]
\scriptsize
\node[end] (source) {Zdrojový \\ kód};
\node (lexer) [right=of source] {Tokenizace \\ (Lexer)};
\node (preparser) [right=of lexer] {Prvotní parsování \\ (Preparser)};
\node (parser) [right=of preparser] {Finální parsování \\ (Parser)};
\node (codegenerator) [right=of parser] {Generace LLVM IR \\ (Code Generator)};
\node (llvminterface) [right=of codegenerator] {Kompilace LLVM IR \\ (LLVM Interface)};
\node[end] (exe) [right=of llvminterface] {Spustitelný \\ soubor};
\draw[arrow] (source) -- (lexer);
\draw[arrow] (lexer) -- (preparser);
\draw[arrow] (preparser) -- (parser);
\draw[arrow] (parser) -- (codegenerator);
\draw[arrow] (codegenerator) -- (llvminterface);
\draw[arrow] (llvminterface) -- (exe);
\end{tikzpicture}
\caption{Jednotlivé fáze compileru}
\end{figure}
\section{Způsob psaní kódu}
Při psaní kódu jsem se snažil být pragmatický. Například jsem nedával příliš velký pozor, zda je všechna alokovaná paměť následně i uvolněna,
protože celý compiler běží jen velmi krátkou dobu a všechna paměť je stejně na konci programu uvolněna. Také jsem byl velmi konzervativní s
používáním standardní knihovny templatů pro C++, (STL) protože mám zkušenost, že přílišné používání STL může vést k nečitelným chybám a
zavádějící logice. Znamenalo to ale, že jsem nevyužil některých užitečných datových struktur, které jazyk C++ nabízí, nebo jsem si je musel
implementovat sám. Kód také nekontroluje, zda je uživatelský vstup správný, a jediný kód na nahlášení chyb slouží k vystopování, která
část compileru chybu našla, místo aby dával přehlednou zprávu o tom, co se pokazilo. To mi umožnilo na kódu rychleji pracovat a dovolilo mi
to dokonce některé fáze kompilace, jako třeba Generátor LLVM kódu několikrát celé přepsat, než jsem byl spokojený s výsledkem.

\section{Kostra}
Kód, kterému říkám \enquote{Kostra} řídí běh celého compileru. Při jeho návrhu jsem napřed napsal kód hlavního souboru. Snažil jsem se,
aby kód hlavního souboru byl co nejčitelnější. Rozložil jsem to tak, že většina obtížné práce je schovaná za příjemně vypadající metodou
\texttt{module.build()}. Na druhou stranu jsem věnoval většinu hlavního souboru uživatelskému rozhraní. Místo nesrozumitelného souboru
na stovky řádků jsem tedy dostal jednoduchý soubor, pomocí kterého může případný vývojář najít, která část projektu je zodpovědná za jaký
úkol. Teprve když byl hlavní soubor hotov, jsem implementoval třídy a metody, které jsem si v hlavním souboru vymyslel. Můj přístup
znamenal, že jsem věděl přesně, co musím implementovat. Neztrácel jsem tedy čas psaním zbytečného kódu.
\subsection{Hlavní soubor}
Níže naleznete celý hlavní soubor. Importuje tři knihovny. První, \texttt{module.hpp} je má implementace pomocných tříd, které mi dávají
přehledné rozhraní pro interakci s projektem. Pod tím je množství řádků, které definují argumenty příkazového řádku, zakončené importem
knihovny \texttt{easyargs.h}. Výhodou \texttt{easyargs.h} je, že díky chytrému využití preprocesoru jazyka C umí měnit svůj obsah a
může tedy vrátit strukturu typu \texttt{args\_t}, která obsahuje zadanou hodnotu pro každý argument. Poslední knihovnou je \texttt{iostream},
což je vestavěná knihovna C++ pro výpis do konzole. Pokud nejste seznámeni s programováním v jazyce C++, nebuďte překvapeni syntaxi
\texttt{std::cout}, je to způsob, jakým se v C++ píše do konzole. Také je nutné podotknout, že knihovna Easyargs je určena pro C, nikoli C++,
tudíž jsem musel řešit rozdíly mezi oběma jazyky.

\begin{listing}[H]
\caption{Celý hlavní soubor}
\inputminted[breaklines]{cpp}{../cppcompiler/src/_main.cpp}
\end{listing}

\subsection{Třída \texttt{Module} a další pomocné třídy}
Objekt třídy \texttt{Module} reprezentuje program, který se nachází ve fázích kompilace mezi tokenizací a generací LLVM IR.
Konstruktor \texttt{Module} žádá pouze soubor s kódem, na kterém jsou nejdříve provedeny první dvě fáze kompilace, v kódu
přezdívané Lexer a Preparser. \texttt{Module} se skládá ze dvou částí, hlavní tabulky symbolů, kterou konstruktor dostane z Preparseru,
a balíčku objektů knihovny LLVM, ve kterém se budou ukládat generované LLVM instrukce. Třída \texttt{Module} má jedinou metodu a to
je \texttt{build}. \texttt{build} se podívá do tabulky symbolů a zkompiluje vše potřebné do LLVM IR. Jako návratovou hodnotu má
\texttt{build} objekt třídy \texttt{IrFile}. Ta je jednou z několika tříd, které mi usnadňují práci se soubory. Přestože se třída
jmenuje \texttt{IrFile}, nekoresponduje se souborem na disku. Tyto pomocné třídy totiž reprezentují jen dočasnou hodnotu, dokud není
zavolána jejich \texttt{commit} metoda, která jim teprve přiřadí trvalé místo na disku.
\subsection{Symboly a Tabulky symbolů}
Základní jednotka kódu se v mém jazyce nazývá Symbol. Jednotlivé symboly v mém compileru reprezentují instance třídy \texttt{Symbol}.
Symbolům jsou přiřazeny tabulky symbolů a ke každému Symbolu je možné přiřadit pouze jednu tabulku. Tabulka symbolů může také vlastním
symbolům přiřazovat jména. Toho využívám v implementaci proměnných a funkcí.

Existuje několik tabulek zároveň. Jednotlivé tabulky mají mezi sebou hierarchickou strukturu. Hlavní tabulka popisuje globálně viditelné
symboly, tedy globální proměnné a funkce, zatímco podřízené tabulky popisují lokální symboly.

Každý symbol v sobě ukládá reference na struktury reprezentující ten stejný symbol v určité fázi kompilace. Díky této struktuře nemusím
překládat celý program z jedné reprezentace na jinou naráz, ale můžu zkompilovat jen ty, které potřebuji. Můžu například zkompilovat
pouze globální funkci \texttt{main} a compiler automaticky určí, které symboly zkompilovat.
\section{Tokenizace (Lexer)}
Tokenizace, neboli lexikální analýza probíhá v mém compileru ve fázi, kterou jsem nazval Lexer. Její podstatou je ze vstupního ještě
nezpracovaného kódu vytvořit seznam elementárních jednotek syntaxu. Syntax mého jazyka se totiž neskládá výhradně ze znaků, nýbrž často
z celých slov. Je proto užitečné být schopný číst kód slovo po slovu, místo písmeno po písmenu. Takovým \enquote{slovům} se říká tokeny.
\footnote{Tokeny a slova nejsou ve všech kontextech zaměnitelné. Například při programování LLM se dlouhá slova většinou skládají z několika tokenů.}
Programovací jazyky často vedle textu do tokenů ukládají další informace. Můj jazyk tohle nedělá, proto (mimo jiné) můj compiler nedokáže
uživateli říct, na jakém místě v souboru našel syntaktickou chybu, protože ta informace je první věcí, které se zbavil. Samozřejmě
to v budoucnosti lze změnit, jenom jsem tenhle nedostatek zpozoroval příliš pozdě a před maturitní zkouškou jsem měl větší problémy.
Lexer byla první část, kterou jsem na projektu dokončil. Od té doby jsem nemusel na Lexeru téměř nic měnit, protože již bohatě stačil
pro mé potřeby. Podstatou je, že se Lexer snaží seskupovat alfanumerické znaky. Pokud Lexer narazí na whitespace, tak ho vynechá, jenom
v tom místě nesmí seskupovat znaky. Speciální znaky mají každý svá vlastní pravidla tokenizace. Pokud je uživatel rozumný, neměl by mít Lexerem
v jazyce potíže.

\begin{figure}[H]
\centering
\begin{minipage}{0.4\linewidth}
\centering
\textbf{Zdrojový kód, jak jej vidí uživatel}
\medskip
\large
\begin{verbatim}
fn num(n) : {
    if (n >= 10) {
        num(n / 10);
    };
    putchar(48+n%10);
};
\end{verbatim}
\vspace{6mm}
\end{minipage}
\begin{minipage}{\linewidth}
\centering
\textbf{Zdrojový kód před tokenizací}
\medskip
\begin{tikzpicture}
\matrix [matrix of nodes,
    nodes={draw, minimum width=8mm, minimum height=8mm,
           text height=1.5ex, text depth=.25ex},
    column sep=0.5mm, row sep=0.5mm] {
f & n & \phantom{X} & n & u & m & ( & n & ) & \phantom{X} & : & \phantom{X} & \{ & \phantom{X} \\
\phantom{X} & i & f & \phantom{X} & ( & n & \phantom{X} & > & = & \phantom{X} & 1 & 0 & ) & \phantom{X} & \{ & \phantom{X}  \\
\phantom{X} & \phantom{X} & n & u & m & ( & n & \phantom{X} & / & \phantom{X} & 1 & 0 & ) & ; & \phantom{X}  \\
\phantom{X} & \} & ; & \phantom{X} \\
\phantom{X} & p & u & t & c & h & a & r & ( & 4 & 8 & + & n & \% & 1 & 0 & ) & ; & \phantom{X}  \\
\} & ; & \phantom{X}  \\
};
\end{tikzpicture}
\vspace{6mm}
\end{minipage}
\begin{minipage}{\linewidth}
\centering
\textbf{Zdrojový kód po tokenizaci}
\medskip
\begin{tikzpicture}
\matrix [matrix of nodes,
    nodes={draw, minimum width=16.5mm, minimum height=5mm,
           text height=1.5ex, text depth=.25ex},
    column sep=1mm, row sep=1mm] {
fn & num & ( & n & ) & : & \{ \\
if & ( & n & >= & 10 & ) & \{ \\
num & ( & n & / & 10 & ) & ; \\
\} & ; \\
putchar & ( & 48 & + & n & \% & 10 & ) & ; \\
\} & ; \\
};
\end{tikzpicture}
\end{minipage}
\caption{Rozdíl kódu před a po tokenizaci}
\end{figure}
\section{První fáze parsování (Preparser)}
Parsování neboli syntaktická analýza rozděluje program na jednotlivé mezi sebou propojené operace. V této části je nutné poznamenat,
že není vůbec běžné rozdělovat parsování na dva kroky a Preparser není obecně zavedený pojem. Parsování jsem rozdělil na dva kroky,
protože je pro moje budoucí plány důležité, abych měl co nejdříve představu o názvech symbolů. Proto jsem vytvořil tuto fázi, která
sestaví tabulku symbolů hned jako druhý krok. Aby dokázala zjistit věci, jako jsou třeba názvy proměnných je potřeba analyzovat
syntax klíčových slov. Aby se nemusela opakovat práce, uloží si tato fáze syntax libovolného klíčového slova ve strukturovaném formátu.
Aby dále usnadnil dalším fázím práci, přidá Preparser nově vytvořeným symbolům značky podle toho jestli se jedná o operátor a případně
o jaký druh operátoru se jedná. Preparser má 4 známé druhy symbolů: hodnota, prefix operátor, infix operátor, postfix operátor. Prefix
operátor jsou operátory, které operují jenom na hodnotě za nimi. Příkladem prefix operátoru je negace ($!x$), nebo předpona minus ($-x$).
Postfix operátory naopak operují na hodnotě před nimi. V programování jsou velmi vzácné. Příkladem by byl post-decrement operátor ($x--$).
Nejčastějším druhem operátoru jsou infix operátory, které operují na obou sousedních hodnotách. Příkladem je odčítání ($x-y$). Pokud
se jedná o složitější operátor, jako jsou například závorky, nedokážeme jej pomocí těchto čtyř druhů samotných popsat. Musíme se
tedy spolehnout na kód pro řešení klíčových slov, který naopak dokáže popsat téměř jakýkoli syntax. Proto jsou závorky klíčovým slovem.
Všechny hodnoty musí být odděleny infix operátory. Aby to Preparser zajistil, vloží prázdné infix operátory tam, kde je to potřeba.

Níže vidíte stejný program jako předtím, který však již prošel Preparserem. Není to optimální příklad, protože obsahuje pouze hodnoty a infix operátory.
Také nevím, jestli do maturity přidám závorky jako postfix operátor pro volání funkce, nebo to bude opravdu vypadat jako na obrázku.
Nicméně graf dole ukazuje, jak vypadá rozložení paměti v Preparseru. Šipky značí, že Symbol, odkud ta šipka vede, má výraz, na který šipka
ukazuje, uložený v sobě a má nad ním absolutní kontrolu. Bílé prázdné čtverečky značí prázdný výraz. (0 symbolů)

\begin{figure}[H]
\centering
\begin{tikzpicture}[
value/.style={rectangle, draw=blue!60, fill=blue!5, very thick, minimum size=8mm},
infix/.style={rectangle, draw=red!60, fill=red!5, very thick, minimum size=8mm},
empty/.style={rectangle, draw=black!60, fill=white!5, very thick, minimum size=8mm, dotted},
]
\node[value] (fn) {fn};
\node[value] (curly1) [above=of fn] {\{};
\node[value] (if) [above=of curly1] {if};
\matrix (m1) [above=of if, matrix of nodes,
              nodes={draw, minimum size=6mm, anchor=center}] {
  |[value]| n & |[infix]| >= & |[value]| 10 \\
};
\node[value] (curly2) [right=of m1] {\{};
\matrix (m2) [above=of curly2, matrix of nodes,
              nodes={draw, minimum size=6mm, anchor=center}] {
  |[value]| num & |[infix]| \phantom{X} & |[value]| ( \\
};
\matrix (m3) [above=of m2-1-3, matrix of nodes,
              nodes={draw, minimum size=6mm, anchor=center}] {
  |[value]| n & |[infix]| / & |[value]| 10 \\
};
\node[empty] (empty1) [right=of m2] {\phantom{X}};
\node[empty] (empty2) [below=of empty1] {\phantom{X}};
\matrix (m4) [right=of empty2, matrix of nodes,
              nodes={draw, minimum size=6mm, anchor=center}] {
  |[value]| putchar & |[infix]| \phantom{X} & |[value]| ( \\
};
\matrix (m5) [above=of m4-1-3, matrix of nodes,
              nodes={draw, minimum size=6mm, anchor=center}] {
  |[value]| 48 & |[infix]| + & |[value]| n & |[infix]| \% & |[value]| 10 \\
};
\node[empty] (empty3) [below=of m4] {\phantom{X}};
\draw[->, very thick] (fn.north) -- (curly1.south);
\draw[->, very thick] (curly1.north) -- (if.south);
\draw[->, very thick] (if.north) -- (m1.south);
\draw[->, very thick, bend right=25] (if.east) to (curly2.south);
\draw[->, very thick] (curly2.north) to (m2);
\draw[->, very thick, bend left=10] (curly2.east) to (empty1.west);
\draw[->, very thick, bend right=25] (if.east) to (empty2.west);
\draw[->, very thick] (m2-1-3) to (m3);
\draw[->, very thick, bend right=25] (curly1.east) to (m4);
\draw[->, very thick] (m4-1-3) to (m5);
\draw[->, very thick, bend right=20] (curly1.east) to (empty3);

\end{tikzpicture}

\caption{Program z předchozí části očima Preparseru}
\end{figure}
\pagebreak
\section{Druhá fáze parsování (Parser)}
V této fázi konečně sestavíme AST. Jedná se o strom uzlů, kde každý uzel znázorňuje nějakou operaci. Díky pilné práci předchozí fáze je
sestavení AST téměř bezbolestné. Bohužel je tu ale jeden problém. Pokud se podíváte na obrázek níže, můžete si všimnout, že nemá stejný
počet symbolů, jako ten nahoře. Naštěstí jsem si byl toho problému vědom a vyřešil jsem ho. Symboly mají chytrou funkci, díky které se
mohou spojovat dohromady, jako například u závorek, které v AST nejsou, nebo je také možné nové symboly přidávat. To je sice trochu
nebezpečnější, protože je teoreticky možné, že dřívější fáze si zažádá symbol vytvořený budoucí fází a ten symbol bude mít jen reprezentaci
z té pozdější fáze, ale zatím v mém limitovaném testování tato situace nenastala, takže je logické usoudit, že u správně napsaného klíčového
slova to nebude problém.
\begin{figure}[H]
\centering
\begin{tikzpicture}[
node/.style={circle, draw=black!60, fill=white!5, very thick, minimum size=4mm},
]
\node[node] (fn) {fn};
\node[node] (curly1) [above=of fn] {\{};
\matrix (m1) [above=of curly1, matrix of nodes,
              nodes={draw, minimum size=6mm, anchor=center}, column sep = 10mm] {
  |[node]| if & |[node]| \phantom{X} & |[node]| 0 \\
};

\matrix (m2) [above=of m1, matrix of nodes,
              nodes={draw, minimum size=6mm, anchor=center}, column sep = 10mm] {
  |[node]| >= & |[node]| \{ & |[node]| 0 & |[node]| putchar & |[node]| + \\
};
\matrix (m3) [above=of m2, matrix of nodes,
              nodes={draw, minimum size=6mm, anchor=center}, column sep = 10mm] {
  |[node]| n & |[node]| 10 & |[node]| \phantom{X} & |[node]| 0 & |[node]| 48 & |[node]| \% \\
};
\matrix (m4) [above=of m3-1-3, matrix of nodes,
              nodes={draw, minimum size=6mm, anchor=center}, column sep = 10mm] {
  |[node]| num & |[node]| / \\
};
\matrix (m5) [above=of m3-1-6, matrix of nodes,
              nodes={draw, minimum size=6mm, anchor=center}, column sep = 10mm] {
  |[node]| n & |[node]| 10\\
};
\matrix (m6) [above=of m4-1-2, matrix of nodes,
              nodes={draw, minimum size=6mm, anchor=center}, column sep = 10mm] {
  |[node]| n & |[node]| 10\\
};
\draw[->, very thick] (fn) -- (curly1);
\draw[->, very thick] (curly1) -- (m1-1-1);
\draw[->, very thick] (curly1) -- (m1-1-2);
\draw[->, very thick] (curly1) -- (m1-1-3);
\draw[->, very thick] (m1-1-1) -- (m2-1-1);
\draw[->, very thick] (m1-1-1) -- (m2-1-2);
\draw[->, very thick] (m1-1-1) -- (m2-1-3);
\draw[->, very thick] (m1-1-2) -- (m2-1-4);
\draw[->, very thick] (m1-1-2) -- (m2-1-5);
\draw[->, very thick] (m2-1-1) -- (m3-1-1);
\draw[->, very thick] (m2-1-1) -- (m3-1-2);
\draw[->, very thick] (m2-1-2) -- (m3-1-3);
\draw[->, very thick] (m2-1-2) -- (m3-1-4);
\draw[->, very thick] (m2-1-5) -- (m3-1-5);
\draw[->, very thick] (m2-1-5) -- (m3-1-6);
\draw[->, very thick] (m3-1-3) -- (m4-1-1);
\draw[->, very thick] (m3-1-3) -- (m4-1-2);
\draw[->, very thick] (m3-1-6) -- (m5-1-1);
\draw[->, very thick] (m3-1-6) -- (m5-1-2);
\draw[->, very thick] (m4-1-2) -- (m6-1-1);
\draw[->, very thick] (m4-1-2) -- (m6-1-2);
\end{tikzpicture}

\caption{AST pro stejný program, jako předtím}
\end{figure}
\pagebreak
\section{Generace LLVM IR (Code Generator)}
\begin{listing}[H]
\begin{minted}{LLVM}
define i64 @num(i64 %0) {
  %2 = icmp sge i64 %0, 10
  %3 = zext i1 %2 to i64
  %4 = icmp ne i64 %3, 0
  br i1 %4, label %5, label %8

5:                                                ; preds = %1
  %6 = sdiv i64 %0, 10
  %7 = call i64 @num(i64 %6)
  br label %9

8:                                                ; preds = %1
  br label %9

9:                                                ; preds = %8, %5
  %10 = phi i64 [ 0, %5 ], [ 0, %8 ]
  %11 = srem i64 %0, 10
  %12 = add i64 48, %11
  %13 = call i64 @putchar(i64 %12)
  ret i64 0
}
\end{minted}
\caption{Funkce num(n) zkompilovaná do LLVM IR}
\end{listing}
Na rozdíl od parsování nemám pro Generaci LLVM obrázek, který by vše podrobně vysvětlil. Na rozdíl od předchozích fází, kde jsem si mohl vše
volně rozhodovat sám, se zde totiž mísí moje názory o strukturování kódu s názory LLVM. Proto byla tato fáze ze všech nejkritičtější a musel
jsem ji třikrát přepsat. Dosud totiž mířily všechny šipky pouze jedním směrem, tudíž symboly vůbec neměly ponětí o tom, na jaké větvi se nachází.
Pro generaci instrukcí však symbol nepotřebuje znát pouze, na jakých hodnotách operuje, ale i do jaké funkce má svou instrukci napsat, takže musí
být schopný se podívat oběma směry. Proto jsem rozdělil tuto fázi na dva kroky.

V prvním kroku procházíme graf a každému navštívenému symbolu řekneme, odkud jsme k němu došli. V prvním kroku zároveň vytváříme proměnné
a funkce. Teprve, když už každý symbol ví, co je v obou směrech, mohou symboly doputovat tím řetězem k funkci, ve které se nacházejí a přidat
do ní patřičnou instrukci.

Kromě toho komplikovaného procesu si schopný programátor v kódu výše něčeho všimne. Program nahoře je totiž velmi špatně napsaný. Hned jako první
se v něm porovnávají dvě čísla za vzniku logické hodnoty, která je pak zkonvertovaná na číslo, a pak zase na logickou hodnotu. Dále funkce skáče
na blok kódu \texttt{\%8}, a pak z něj ihned skočí na blok kódu \texttt{\%9}. Pokud se podíváte na instrukci \texttt{phi}, je očividné, že jde
zjednodušit pouze na \texttt{0}, a že její vypočtená hodnota není nikdy použita.

Jak tedy vidíte můj compiler zatím při generaci LLVM IR nedělá vůbec žádnou optimalizaci. On vlastně nedělá optimalizaci nikde. To je v souladu
s mými cíli, protože by implementovat pořádné optimalizace ubíralo času od implementace mnohem důležitějších věcí.

\section{Rozhraní pro LLVM (LLVM Interface)}
Tohle je poslední fáze. Zajímavá je v tom, že jako jediná probíhá zcela mimo pomocnou třídu \texttt{Module}. Je také, alespoň dle mého názoru,
nejméně zajímavá. Jediné, co dělá, je že volá knihovnu LLVM, aby zkompilovala daný LLVM IR kód do Assembly, nebo objektového kódu. Případně
pak zavolá linker, (já si vybral LLD) který pak ze souboru s objektovým kódem udělá spustitelný soubor. Kromě toho nemá smysl zacházet tady do větších detailů, protože jsem velkou část zkopíroval z tutoriálu na stránkách LLVM, takže pokud vás nejen tato část projektu zaujala, zkuste si to přečíst:
\href{https://llvm.org/docs/tutorial/MyFirstLanguageFrontend/index.html}{https://llvm.org/docs/tutorial/MyFirstLanguageFrontend/index.html}

\section{Problémy s Windows}
Následující část obsahuje primárně mé subjektivní názory a mou osobní zkušenost. Není nijak důležitá k pochopení vnitřní struktury mého
programovacího jazyka. Důležité je říct, že moje negativní zkušenost s Windows nebyla způsobená mým primárním zaměřením na Linux,
i když skutečnost, že nemám neomezený přístup k počítači s Windows, která mě přinutila používat zamrzávající VM, situaci nepomohla.
Následuje převyprávění několikaměsíční snahy zkompilovat můj projekt na školním počítači.

\subsection{První pokus -- Visual Studio}
Ze začátku jsem se pokoušel zprovoznit kompilaci přes Visual Studio, kanonický nástroj pro vytváření Windows native aplikací. Prvním problémem
bylo, že nemohu prostě vložit svůj Makefile do Visual Studia a očekávat, že to bude fungovat. Visual Studio podporuje tři build systémy pro C++:
MSBuild, CMake a NMake. Jeden z nich jsem si musel vybrat. MSbuild jsem ihned zamítl, protože funguje na základě složitě navigovatelných menu,
takže to bylo mezi CMake a NMake. S CMake mám předchozí zkušenost, tedy hlavně negativní předchozí zkušenost, zatímco NMake je port Makefile
pro Windows. Bohužel ale Makefile a NMake jsou dostatečně odlišné, že bych stejně musel mít soubory na kompilaci dva. Protože jsem neměl náladu
ani energii, rozhodl jsem se pro CMake. Myslel jsem si, že umělá inteligence, která je trénovaná na mnohem populárnějším CMake mi dá, alespoň
v prvních pár pokusech, alespoň částečně funkční build skript. Ukázalo se ale, že Visual Studio nemá žádný bezbolestný způsob, jak si stáhnout
knihovnu LLVM, takže jsem si ji musel stáhnout manuálně, nebo si ji dokonce zkompilovat sám. Nemyslím si, že projekt rozměru mého má právo žádat
o uživateli zkompilovat knihovnu velikosti LLVM, takže jsem si začal stahovat LLVM, když jsem dostal nápad, čímž tento pokus oficiálně skončil.

\subsection{Druhý pokus -- Terminál}
Můj nápad byl jednoduchý: \enquote{Jestli si stejně musím stahovat celé LLVM, co mi tedy poskytuje Visual Studio?} Uvědomil jsem si totiž,
že mimo jiné distribuce LLVM obsahuje Clang, compiler pro C a C++. Jestliže mám tedy svůj kód a C++ compiler, nic mi nebrání prostě Visual Studio
přeskočit a rovnou spouštět compiler pomocí terminálu. Co mi poskytuje build systém, který mě stejně odpuzuje, když kompilační skript můžu napsat
v jakémkoli skriptovacím jazyce schopném spouštět příkazy přes příkazový řádek. Takže jsem nechal ChatGPT mi vygenerovat BATCH skript, který
prostě spustí všechny potřebné kompilační příkazy. Tohle fungovalo.

Důraz zde dávám na minulý čas. Chtěl jsem přinést ukázat svůj projekt na školním počítači. Když jsme však přišli do učebny, zhrozili jsme se.
Školní administrace totiž v ten stejný týden (možná dokonce ten předchozí den) aktualizovala počítače v učebně na Windows 11. Když jsme se
přihlásili ukázalo se, že upgrade nevyšel úplně hladce; počítače každých několik minut ukázaly modrou obrazovku. To však nebyl můj největší
problém. Ten byl v tom, že upgrade odinstaloval Visual Studio.

Proč mi vadilo, že nemám Visual Studio, když jsem ho neplánoval použít? Protože Visual Studio, jako kononický nástroj pro vytváření Windows
native aplikací není pouze editorem. Je to správce celé řady nástrojů, které jsou potřeba při vývoji softwaru. Pokud bylo Visual Studio pryč,
byly pryč i vývojářské nástroje jazyka C++. Pro zkompilování C++ nestačí, jak jsem si myslel, C++ kód a C++ compiler, ale je potřeba i
C++ standardní knihovny, kterou jsem kvůli absenci Visual Studia neměl a k instalaci Visual Studia je třeba administrátorského přístupu.
Radši, než vysvětlovat, že potřebuji Visual Studio nainstalované přestože ho neplánuji používat, jsem se rozhodl najít alternativní metodu.

\subsection{Třetí pokus -- MinGW}
Navenek vypadá MinGW jako nástroj, který jsem měl použít rovnou. Poskytuje na Windows prostředí unixového příkazového řádku, což by mi
umožnilo kompilovat Linux i Windows verze jedním Makefile. Navíc k instalaci nepotřebuje administrátorská práva. Má však i nevýhody.
Hlavní nevýhodou z mého hlediska je, že výsledný soubor je závislý na kódu uvnitř MinGW, (zvláště pokud je to zkompilované s MSYS2 runtime)
takže jediný způsob, jak ten projekt spustit mimo prostředí MinGW je najít všechny .dll soubory, na kterých zkompilovaný soubor závisí a
přesunout je do stejné složky, což není něco, co by měl člověk dělat. Doufám tedy, že pokud jste to dočetli až sem, tak mi prominete, že
můj projekt pravděpodobně MinGW na Windows neopustí.

\subsection{Změny v kódu}
V některých částech kódu, hlavně v Code Generatoru a LLVM interface, se vyskytuje kód, který by na Windows nefungoval. Používám proto
C preprocesoru, abych mohl poskytnout oběma platformám funkční implementaci. Jeden problém, který se mi vyskytl pouze na Windows je,
že C++ knihovna, která mi poskytovala rozhraní mezi mým compilerem a LLD, byla rozbitá a nefungovala, takže jsem specificky pro Windows
udělal alternativní implementaci, která prostě volá LLD přes příkazový řádek.

Větším problémem je, že Linux a Windows používají pro komunikaci s operačním systémem velmi odlišné rozhraní. Kvůli tomu jsem musel udělat
statickou knihovnu, která dělá jednu věc: importuje 3 Windows API funkce. Program přeložený do objektového kódu pak s knihovnou zkombinuji
pomocí linkeru. Na linuxu mi stačila \texttt{syscall} instrukce v Assembly.

\subsection{Nedostatky Windows verze}
Kromě závislosti na MinGW je momentálně největší problém v zobrazování Unicode znaků na obrazovce. Problém to je v C/C++ na Windows dlouhodobý.
Windows očekává, že výstup obsahující Unicode znaky mu bude předán v UTF-16, zatímco Linux přijímá UTF-8 bez problémů. Easyargs je cíleno na unixové
prostředí, takže kód očekává 8-bitové znaky, jako v UTF-8. Když se Easyargs snaží předat Windowsu text v UTF-8, Windows odmítne rozpoznat Unicode
znaky a napíše místo toho několik náhodně vypadajících znaků.

\chapter{Uživatelská dokumentace}
\section{Instalace a spuštění}
Následuje vysvětlení postupu kompilace pro Windows přes MinGW. Aktuálně není na GitHubu projekt ve zkompilované verzi, proto je nutné si jej
zkompilovat sám.

\subsection{Kompilace na Linuxu}
Následuje vysvětlení postupu kompilace pro Linux. Aktuálně není na GitHubu projekt ve zkompilované verzi, proto je nutné si jej zkompilovat sám.

\paragraph{Předpoklady:}
\begin{itemize}
 \item Git (pouze ke stažení projektu z GitHubu)
 \item Funkční C++ compiler (clang++ nebo g++)
 \item Make
 \item LLVM
 \item LLD
 \item Příkazový řádek
\end{itemize}
\paragraph{Krok 1:}
Nejprve je důležité mít kód stažený. Zde použijeme rozhraní Gitu pro příkazový řádek, ale je také možné využít různých grafických rozhraní, která
dokáží udělat to stejné. Alternativou ke Gitu by bylo si projekt z GitHubu stáhnout jako .zip archiv. Otevřeme si tedy náš příkazový řádek a ke stažení si projektu použijeme příkaz:
\begin{verbatim}
git clone https://github.com/F100cTomas/Lang.git
\end{verbatim}
Stažený kód by se nyní měl nacházet ve složce Lang/. Pro další kroky budeme chtít přejít do složky s Makefile souborem:
\begin{verbatim}
cd Lang/cppcompiler/
\end{verbatim}
\paragraph{Krok 2:}
Nyní ze složky \texttt{cppcompiler/} spustíme \texttt{make}. Příkazu \texttt{make} je možné poskytovat argumenty. V aktuální verzi je možné
poskytnout pouze dva. První je \texttt{BUILD}, kterým specifikujeme, že nechceme, aby výsledný soubor obsahoval debugovací symboly.
Druhý je \texttt{CXX}, který říká, jaký C++ compiler bude použit. Protože výchozí hodnota pro \texttt{BUILD} je \texttt{debug}, doporučuji
specifikovat hodnotu \texttt{release}. Pro rychlejší kompilaci použijeme argument \texttt{-j8}, který práci rozdělí mezi více vláken. Příkaz
vypadá takhle:
\begin{verbatim}
make -j8 BUILD=release CXX=clang++
\end{verbatim}
\paragraph{Krok 3:}
Vzniklý soubor se jmenuje \texttt{lang}. Pro spuštění compileru použijeme:
\begin{verbatim}
./lang code
\end{verbatim}
Pro vypsání tabulky argumentů \texttt{lang} existuje:
\begin{verbatim}
./lang code -h
\end{verbatim}
Pokud chcete, aby byl příkaz \texttt{lang} přístupný všude, můžete poté soubor přesunout do složky \texttt{/usr/bin/}, nebo jiné složky v
proměnné prostředí \texttt{PATH}, nebo složku s compilerem do \texttt{PATH} přidat.

\subsection{Kompilace na Windows}
Kompilace na Windows je téměř úplně stejná, jak na Linuxu. Stačí pouze nahradit Linux příkazový řádek příkazovým řádkem MinGW.
Důležité je mít všechen předpokládaný software stažený uvnitř prostředí MinGW.

\section{Uživatelské rozhraní}
Uživatel s projektem interaguje pomocí příkazového řádku. Jako první argument uvádí vždy uživatel soubor, který bude zkompilován.
Dále může uživatel specifikovat název vzniklého souboru pomocí možnosti \texttt{-o <název souboru>}. \texttt{-a} způsobí, že
místo spustitelného souboru uživatel dostane instrukce procesoru ve formátu Assembly. \texttt{-d} vypíše do konzole mnoho nezajímavého
textu, který slouží hlavně pro debugování compileru, nikoli uživatelských programů.
\begin{figure}[H]
\begin{verbatim}
USAGE:
    ./lang <vstup> [OPTIONS]

ARGUMENTS:
    <vstup>         Soubor ke kompilaci

OPTIONS:
    -o <výstup>    Název zkompilovaného souboru (default: program)
    -a              Assembly výstup
    -d              Debug výstup
    -h              Seznam argumentů
\end{verbatim}
\caption{Nápověda, kterou uživatelské rozhraní poskytuje}
\end{figure}

\section{Funkce}
Každý program začíná ve funkci \texttt{main()}. Návratová hodnota \texttt{main()} koresponduje s kódem ukončení programu.
Tohle je nejjednodušší možný program, hned se ukončí s kódem 0.
\begin{verbatim}
fn main() : 0;
\end{verbatim}
Pro tisknutí na obrazovku stačí zadat funkci \texttt{putchar} ASCII hodnotu požadovaného znaku. Následující program tedy vytiskne na obrazovku \enquote{A}:
\vspace{1em} \\
\begin{center}
\begin{minipage}{0.4\linewidth}
\begin{verbatim}
fn main() {
  putchar 65;
  putchar 10;
  0
};
\end{verbatim}
\end{minipage}
\begin{minipage}{0.5\linewidth}
Každý znak má v této tabulce určitou hodnotu, \texttt{putchar} na základě toho ten určitý znak vypíše.
\end{minipage}
\end{center}
\begin{table}[H]
\centering
\begin{tabularx}{\textwidth}{|*{8}{>{\centering\arraybackslash}X|}}
\hline
Číslo & Znak & Číslo & Znak & Číslo & Znak & Číslo & Znak \\ \hline
0  & NUL & 32 & SPACE & 64  & @ & 96  & ` \\ \hline
1  & SOH & 33 & !     & 65  & A & 97  & a \\ \hline
2  & STX & 34 & "     & 66  & B & 98  & b \\ \hline
3  & ETX & 35 & \#    & 67  & C & 99  & c \\ \hline
4  & EOT & 36 & \$    & 68  & D & 100 & d \\ \hline
5  & ENQ & 37 & \%    & 69  & E & 101 & e \\ \hline
6  & ACK & 38 & \&    & 70  & F & 102 & f \\ \hline
7  & BEL & 39 & '     & 71  & G & 103 & g \\ \hline
8  & BS  & 40 & (     & 72  & H & 104 & h \\ \hline
9  & TAB & 41 & )     & 73  & I & 105 & i \\ \hline
10 & LF  & 42 & *     & 74  & J & 106 & j \\ \hline
11 & VT  & 43 & +     & 75  & K & 107 & k \\ \hline
12 & FF  & 44 & ,     & 76  & L & 108 & l \\ \hline
13 & CR  & 45 & -     & 77  & M & 109 & m \\ \hline
14 & SO  & 46 & .     & 78  & N & 110 & n \\ \hline
15 & SI  & 47 & /     & 79  & O & 111 & o \\ \hline
16 & DLE & 48 & 0     & 80  & P & 112 & p \\ \hline
17 & DC1 & 49 & 1     & 81  & Q & 113 & q \\ \hline
18 & DC2 & 50 & 2     & 82  & R & 114 & r \\ \hline
19 & DC3 & 51 & 3     & 83  & S & 115 & s \\ \hline
20 & DC4 & 52 & 4     & 84  & T & 116 & t \\ \hline
21 & NAK & 53 & 5     & 85  & U & 117 & u \\ \hline
22 & SYN & 54 & 6     & 86  & V & 118 & v \\ \hline
23 & ETB & 55 & 7     & 87  & W & 119 & w \\ \hline
24 & CAN & 56 & 8     & 88  & X & 120 & x \\ \hline
25 & EM  & 57 & 9     & 89  & Y & 121 & y \\ \hline
26 & SUB & 58 & :     & 90  & Z & 122 & z \\ \hline
27 & ESC & 59 & ;     & 91  & [ & 123 & \{ \\ \hline
28 & FS  & 60 & <     & 92  & \textbackslash & 124 & | \\ \hline
29 & GS  & 61 & =     & 93  & ] & 125 & \} \\ \hline
30 & RS  & 62 & >     & 94  & \^{} & 126 & \~{} \\ \hline
31 & US  & 63 & ?     & 95  & \_ & 127 & DEL \\ \hline
\end{tabularx}
\caption{ASCII}
\end{table}
\begin{center}
\begin{minipage}{0.4\linewidth}
\begin{verbatim}
fn num(n) : {
    if (n >= 10) {
        num(n / 10);
    };
    putchar(48+n%10);
};
\end{verbatim}
\end{minipage}
\begin{minipage}{0.5\linewidth}
Funkce mohou mít parametr a být rekurzivní.
\end{minipage}
\end{center}
\vspace{1em}
\section{Operátory}
\vspace{1em}
Jazyk podporuje jednoduché aritmetické operátory, ($+,\ -,\ *,\ /,\ \%$) porovnávací operátory, ($==,\ !=,\ >=,\ >,\ <=,\ <$) a přiřazovací operátor. ($=$)
\section{Klíčová slova \texttt{if} a \texttt{while}}
\vspace{1em}
Klíčová slova \texttt{if} i \texttt{while} se chovají, jak by člověk očekával. Možná někoho zarazí, že if vrací hodnotu na základě toho, jestli byla podmínka
pravdivá.
\vspace{1em}
\section{Proměnné}
\vspace{1em}
\begin{center}
\begin{minipage}{0.4\linewidth}
\begin{verbatim}
let global = 66;
fn main() {
  let local = 9;
  global = global - 1;
  local = local + 1;
  putchar global;
  putchar local;
};
\end{verbatim}
\end{minipage}
\begin{minipage}{0.5\linewidth}
Existují lokální i globální proměnné. Jejich hodnoty se mohou měnit.
\end{minipage}
\end{center}
\section{Ukázkové programy}
\begin{center}
\begin{minipage}{0.45\linewidth}
\textbf{FizzBuzz}
\begin{verbatim}
 fn fizz(f) : {
	if (f % 3 == 0) {
		putchar (102);
		putchar (105);
		putchar (122);
		putchar (122);
	} else 0;
};
fn buzz(b) : {
	if (b % 5 == 0) {
		putchar (98);
		putchar (117);
		putchar (122);
		putchar (122);
	} else 0;
};
fn number(num) : {
	if (num >= 10) {
		number (num / 10);
	} else 0;
	putchar (48 + num % 10);
};
fn fizzbuzz1(n1) : {
	if (n1 % 3 != 0) {
		fizzbuzz2 (n1);
	} else 0;
};
fn fizzbuzz2(n2) : {
	if (n2 % 5 != 0) {
		number (n2);
	} else 0;
};
fn main() : {
	let i = 1;
	while (i < 100) {
		fizz (i);
		buzz (i);
		fizzbuzz1 (i);
		putchar (10);
		i = i + 1;
	};
};
\end{verbatim}
\end{minipage}
\begin{minipage}{0.45\linewidth}
\textbf{Fibonacciho posloupnost}
\begin{verbatim}
fn print1() : {
    putchar 32;
    putchar 49;
    putchar 10;
};

fn number(num) : {
    if (num >= 10) {
        number (num / 10);
    } else 0;
    putchar (48 + num % 10);
};

let a1 = 1;
let a2 = 1;

fn main() : {
    let i = 0;
    print1();
    print1();
    while (i < 50) {
        let a3 = a1 + a2;
        putchar 32;
        number a3;
        putchar 10;
        a1 = a2;
        a2 = a3;
        i = i + 1;
    };
};
\end{verbatim}
\end{minipage}
\end{center}

\chapter{Závěr}
S výsledkem svého projektu jsem spokojen. Všech hlavních cílů se mi podařilo včas dosáhnout. Naopak s vedlejšími cíly jsem byl příliš ambiciózní, což jsem i
čekal, protože jsem si spíše dal při výběru vedlejších cílů prostor, abych mohl projekt rozvíjet různými směry. Upřímně jsem trochu doufal, že stihnu i typový
systém, protože typové systémy jsou velmi důležitou součástí jazyků, ale bohužel tomu tak nebylo. Proto bude typový systém mou hlavní prioritou při budoucí
práci na projektu.

Kromě toho je prostoru na rozvoj velké množství. Jedním obrovským tématem, kterého jsem se ani nedotkl, jsou optimalizace. Jiným milníkem, od kterého je
jazyk stále velmi daleko, je OOP. Nejraději bych, aby se jazyk dostal tak daleko, že mohu jeho compiler přepsat v tom samém jazyce, ale to v momentální
době ještě nevypadá realisticky. Tyto cíle však spadají do daleké budoucnosti a krátkodobě se budu zaměřovat spíše na zvládnutí toho, co jsem si napsal do vedlejších
cílů a ještě není hotové.
\end{document}
